% GENERATED FILE. Edit canonical lessons, then run npm run generate:books.
% canonical-source-hash: fnv1a64:8308ea8c1b77f5d0
% canonical-lessons: FR-C90-difficile, FR-C90-cher, FR-C90-haut, FR-C90-plein, FR-C90-vide

\chapter{Difficult, Dear, High, Full, Empty}
\label{ch:fr-difficult-dear-high-full-empty}
\hlchaptermodality{90}

\begin{chapteropening}
\textbf{By the end of this chapter:} \emph{I can name difficult, dear, high, full, empty.}

Complete the last lesson of chapter 90: I can name difficult, dear, high, full, empty.
\end{chapteropening}

\section[difficile]{difficile — difficult}

\label{lesson:FR-C90-difficile}



\begin{quote}
Before the new one: say the French for fast, then the French for easy.
\end{quote}

\subsection*{You'll want to know: difficile}
\textbf{difficile} — \textquotedblleft{}difficult\textquotedblright{}.

\begin{cousinweb}
\textbf{difficile} $\leftarrow$ Latin \emph{difficilis}, \textquotedblleft{}hard to do\textquotedblright{}: \emph{facile} with \emph{dis-}, \textquotedblleft{}not\textquotedblright{}, in front — a learned borrowing. (EtymDB 2.1)
\end{cousinweb}

\begin{grammarlens}[title={What you've built}]
One more word for this chapter.
\end{grammarlens}

\subsection*{Your turn}
\begin{itemize}
\raggedright
  \item \emph{Say it:} \emph{difficile}
  \item \emph{Say it:} \emph{difficile}, once more
  \item \emph{Say it:} say \emph{difficile}
  \item \emph{Recall:} say the French for fast, then the French for easy, then say \emph{difficile} again
\end{itemize}

\begin{tcolorbox}[breakable,colback=teal!4,colframe=teal!35!black,title={Before you move on}]
What does difficile mean? (\textquotedblleft{}Difficult\textquotedblright{}.) Say it once more.
\end{tcolorbox}
\section[cher]{cher — dear; expensive}

\label{lesson:FR-C90-cher}



\begin{quote}
Before the new one: say the French for easy, then the French for difficult.
\end{quote}

\subsection*{You'll want to know: cher}
\textbf{cher} — \textquotedblleft{}dear; expensive\textquotedblright{}.

\begin{cousinweb}
\textbf{cher} $\leftarrow$ Latin \emph{carus}, \textquotedblleft{}dear, beloved, costly\textquotedblright{} (EtymDB 2.1) — inherited: Latin \emph{ca-} became \emph{che-}.
\end{cousinweb}

\begin{grammarlens}[title={What you've built}]
One more word for this chapter.
\end{grammarlens}

\subsection*{Your turn}
\begin{itemize}
\raggedright
  \item \emph{Say it:} \emph{cher}
  \item \emph{Say it:} \emph{cher}, once more
  \item \emph{Say it:} say \emph{cher}
  \item \emph{Recall:} say the French for easy, then the French for difficult, then say \emph{cher} again
\end{itemize}

\begin{tcolorbox}[breakable,colback=teal!4,colframe=teal!35!black,title={Before you move on}]
What does cher mean? (\textquotedblleft{}Dear; expensive\textquotedblright{}.) Say it once more.
\end{tcolorbox}
\section[haut]{haut — high, tall}

\label{lesson:FR-C90-haut}



\begin{quote}
Before the new one: say the French for difficult, then the French for dear; expensive.
\end{quote}

\subsection*{You'll want to know: haut}
\textbf{haut} — \textquotedblleft{}high, tall\textquotedblright{}.

\begin{cousinweb}
\textbf{haut} $\leftarrow$ Latin \emph{altus}, \textquotedblleft{}high\textquotedblright{}, blended with Frankish \emph{*hauh}, \textquotedblleft{}high\textquotedblright{}, which gave it its \emph{h} (EtymDB 2.1 gives the Frankish side) — inherited.
\end{cousinweb}

\begin{grammarlens}[title={What you've built}]
One more word for this chapter.
\end{grammarlens}

\subsection*{Your turn}
\begin{itemize}
\raggedright
  \item \emph{Say it:} \emph{haut}
  \item \emph{Say it:} \emph{haut}, once more
  \item \emph{Say it:} say \emph{haut}
  \item \emph{Recall:} say the French for difficult, then the French for dear; expensive, then say \emph{haut} again
\end{itemize}

\begin{tcolorbox}[breakable,colback=teal!4,colframe=teal!35!black,title={Before you move on}]
What does haut mean? (\textquotedblleft{}High, tall\textquotedblright{}.) Say it once more.
\end{tcolorbox}
\section[plein]{plein — full}

\label{lesson:FR-C90-plein}



\begin{quote}
Before the new one: say the French for dear; expensive, then the French for high.
\end{quote}

\subsection*{You'll want to know: plein}
\textbf{plein} — \textquotedblleft{}full\textquotedblright{}.

\begin{cousinweb}
\textbf{plein} $\leftarrow$ Latin \emph{plenus}, \textquotedblleft{}full\textquotedblright{} — inherited. (EtymDB 2.1)
\end{cousinweb}

\begin{grammarlens}[title={What you've built}]
One more word for this chapter.
\end{grammarlens}

\subsection*{Your turn}
\begin{itemize}
\raggedright
  \item \emph{Say it:} \emph{plein}
  \item \emph{Say it:} \emph{plein}, once more
  \item \emph{Say it:} say \emph{plein}
  \item \emph{Recall:} say the French for dear; expensive, then the French for high, then say \emph{plein} again
\end{itemize}

\begin{tcolorbox}[breakable,colback=teal!4,colframe=teal!35!black,title={Before you move on}]
What does plein mean? (\textquotedblleft{}Full\textquotedblright{}.) Say it once more.
\end{tcolorbox}
\section[vide]{vide — empty}

\label{lesson:FR-C90-vide}



\begin{quote}
Before the new one: say the French for high, then the French for full.
\end{quote}

\subsection*{You'll want to know: vide}
\textbf{vide} — \textquotedblleft{}empty\textquotedblright{}.

\begin{cousinweb}
\textbf{vide} $\leftarrow$ Old French \emph{vuit} $\leftarrow$ Vulgar Latin \emph{*vocitus}, \textquotedblleft{}empty\textquotedblright{} (EtymDB 2.1), related to Latin \emph{vacuus}, \textquotedblleft{}empty\textquotedblright{} (Lewis \& Short) — inherited.
\end{cousinweb}

\begin{grammarlens}[title={What you've built}]
One more word for this chapter.
\end{grammarlens}

\subsection*{Your turn}
\begin{itemize}
\raggedright
  \item \emph{Say it:} \emph{vide}
  \item \emph{Say it:} \emph{vide}, once more
  \item \emph{Say it:} say \emph{vide}
  \item \emph{Recall:} say the French for high, then the French for full, then say \emph{vide} again
\end{itemize}

\begin{tcolorbox}[breakable,colback=teal!4,colframe=teal!35!black,title={Before you move on}]
What does vide mean? (\textquotedblleft{}Empty\textquotedblright{}.) Say it once more.
\end{tcolorbox}
